\chapter{Problemas técnicos y soluciones}
\label{chap:problemas}

Este capítulo reúne los problemas más costosos del desarrollo. Para cada uno se describen el
síntoma, la causa y la solución, tal como quedaron documentados en el código y en el repositorio.
Muchos de ellos no se deben a un error de programación, sino a detalles de los componentes que no
son evidentes en una primera lectura de su documentación.

\section{UART sobre PIO y latencia del enlace Modbus}
\label{sec:prob-uart}

\paragraph{Síntoma.} Tras migrar el firmware del ESP32-S3 a la Pico 2, las órdenes tardaban entre 3 y
\SI{4}{\second} en tener efecto.

\paragraph{Causa.} En el ESP32-S3 la UART es un periférico hardware: los bytes de la respuesta del
esclavo llegan a una FIFO por interrupción, y cuando ModbusMaster los pide ya están disponibles, así
que cada transacción dura 1 a \SI{3}{\milli\second}. En la Pico, el trazado de la placa lleva la
transmisión a GP14 y la recepción a GP13, y se usó \codigo{SerialPIO}, una UART implementada con una
máquina de estado PIO que funciona en cualquier pin
\verificar{confirmar si el RP2350 permite la función UART en GP14 mediante su función alternativa}.
En la práctica, con \codigo{SerialPIO} el bucle de espera de la respuesta de ModbusMaster agotaba
su tiempo límite, de unos \SI{20}{\milli\second}, en muchas transacciones. Diez escrituras
consecutivas bloqueaban así el núcleo \SI{200}{\milli\second}, y al acumular las lecturas de
supervisión se llegaba a varios segundos.

Un segundo detalle afectaba al cambio de sentido del transceptor. Llamar a \codigo{flush()} antes de
desactivar el emisor, como se haría con una UART hardware, rompía la recepción inmediata de la
respuesta con \codigo{SerialPIO}. La solución fue esperar un tiempo fijo de \SI{200}{\micro\second}
tras el último byte y desactivar entonces DE/$\overline{\text{RE}}$.

\paragraph{Primera solución: doble núcleo.} La versión \fichero{firmwarerasp} dedicó el núcleo 1 al bus
Modbus. El núcleo 0 publicaba la consigna en variables \codigo{volatile} y el núcleo 1 la escribía en
bucle, leyendo velocidad y par cada diez escrituras. Para tomar el bus desde el núcleo 0 (por ejemplo,
para configurar el servo), había que desactivar el trabajador y esperar a que su indicador
\codigo{busy} bajara. El núcleo 0 quedaba libre, pero el diseño tenía inconvenientes: variables
compartidas sin atomicidad, un protocolo de relevo con esperas, y lecturas del servo cuyo instante no
estaba ligado al de la muestra de telemetría.

\paragraph{Solución final: dueño único con lazo de periodo fijo.} Una vez ajustados los tiempos de
espera y el cambio de sentido, cada transacción volvió a durar unos pocos milisegundos. La versión
final prescinde del segundo núcleo: un único lazo de \SI{50}{\milli\second} hace exactamente tres
transacciones por vuelta, en posiciones fijas. El resultado es un tiempo de trabajo de
\SI{11.0}{\milli\second} con sólo \SI{0.2}{\milli\second} de dispersión
(sección~\ref{sec:res-temporizacion}), sin condiciones de carrera, y con lecturas, lógica y escritura
en la misma vuelta y con la misma marca de tiempo.

\section{Configuración del ZSC31014: la ventana del modo comando}
\label{sec:prob-zsc}

\paragraph{Contexto.} Para cambiar la ganancia o el cero del acondicionador hay que escribir en su
EEPROM, y para ello hay que entrar en modo comando. Según la hoja de datos \cite{zsc31014}, tras un
reinicio por encendido (POR) el chip espera la orden Start\_CM durante una ventana de \SI{6}{\milli\second}
si la EEPROM no está bloqueada, o de \SI{1.5}{\milli\second} si lo está. Pasada la ventana, entra en
funcionamiento normal e ignora Start\_CM sin dar ninguna señal de error.

\subsection{Primer obstáculo: provocar el reinicio}

\paragraph{Síntoma.} Tras reprogramar la Pico o reiniciarla por USB, la configuración nunca se
aplicaba.

\paragraph{Causa.} La placa Click seguía alimentada: sin un POR no se abre la ventana.

\paragraph{Solución.} Controlar la alimentación de la Click con su pin EN (GPIO1) y hacer un ciclo
completo de apagado y encendido antes de cada intento.

\subsection{Segundo obstáculo: alimentación parásita por el bus}

\paragraph{Síntoma.} Con EN a nivel bajo, el chip seguía respondiendo por I\textsuperscript{2}C. La
tensión de la Click, medida con un polímetro, se quedaba en unos \num{2.5} a \SI{2.7}{\volt} en lugar
de caer a cero.

\paragraph{Causa.} Las resistencias de \emph{pull-up} de SDA, SCL e INT cuelgan del \SI{3.3}{\volt}
permanente del zócalo mikroBUS, no del carril gobernado por EN. Con el chip sin alimentar, la
corriente entraba por esos pines, atravesaba los diodos de protección internos y alimentaba el carril
del chip. La hoja de datos sitúa el umbral de POR entre \num{1.8} y \SI{2.5}{\volt}: con
\num{2.5}--\SI{2.7}{\volt} de alimentación parásita, el chip nunca veía un reinicio.

\paragraph{Diagnóstico.} Para confirmarlo se añadió la orden LC\_HOLD, que mantiene la célula sin
alimentación \SI{8}{\second} sin bloquear el lazo, con dos modos para las líneas del bus: en alta
impedancia o forzadas a nivel bajo. Las medidas mostraron que sólo forzándolas a nivel bajo la
tensión caía a cero. Hubo además un error en el propio diagnóstico: la primera versión de la prueba
del pin de reinicio dejaba las líneas en alto mientras medía, así que siempre ``detectaba''
realimentación y acusaba al pin EN sin motivo. Se corrigió para que midiera en exactamente las mismas
condiciones que el ciclo real.

\paragraph{Solución.} Durante el apagado se libera el periférico I\textsuperscript{2}C y SDA, SCL e
INT se configuran como salidas a nivel bajo (\codigo{LC\_PARK\_DEFAULT = 1}). El bus se restaura
\emph{antes} de volver a dar tensión, porque \codigo{Wire.begin()} tarda lo bastante como para comerse
la ventana.

\subsection{Tercer obstáculo: acertar la ventana}

\paragraph{Síntoma.} Incluso con un POR limpio, Start\_CM no siempre entraba.

\paragraph{Causa.} La orden tiene que llegar \emph{después} de que el chip termine de arrancar y
\emph{antes} de que se cierre la ventana. El tiempo de arranque depende de los condensadores y del
conmutador de alimentación de la placa Click, de modo que los retardos copiados de otros controladores
(por ejemplo, los \SI{500}{\micro\second} de un controlador de mbed) no servían. Además, cualquier
operación entre el encendido y Start\_CM, aunque sea una impresión por el puerto serie, desplaza el
instante de envío.

\paragraph{Solución.} En lugar de adivinar el retardo, se mide:
\begin{enumerate}
  \item tras dar tensión, se sondea la dirección del chip hasta que responde (unos
        \SI{90}{\micro\second} por sondeo), lo que marca el fin del arranque;
  \item se espera un retardo que cambia en cada intento, de \num{0} a \SI{1050}{\micro\second} en
        pasos de \SI{150}{\micro\second}, a lo largo de hasta ocho ciclos de alimentación;
  \item se envía Start\_CM y se confirma la entrada leyendo B\_Config, cuya respuesta debe empezar por
        \codigo{0x5A};
  \item el intento que funciona se devuelve a la interfaz (paquete \codigo{0x04}), para poder fijar el
        retardo más adelante.
\end{enumerate}
Si ningún intento funciona, se clasifica la causa: ``no contesta'' si nunca hubo ACK, ``el reinicio no
corta'' si sigue respondiendo con EN a nivel bajo, o ``no entra en la ventana''.

\subsection{Cuarto obstáculo: integridad de la EEPROM}

La firma MISR que protege la EEPROM se recalcula al salir del modo comando con Start\_NOM, y se
comprueba en el siguiente POR. Si se corta la alimentación antes de que la firma quede grabada, el chip
arranca desde entonces en estado de diagnóstico y la placa queda inservible. Por eso, tras cualquier
escritura se envía Start\_NOM y se espera diez veces el tiempo típico antes de cortar la tensión. Como
ZMDI\_Config1 sólo se carga tras un POR, la verificación se hace en una segunda pasada, desde un
reinicio limpio.

El firmware fuerza siempre el interfaz I\textsuperscript{2}C en ZMDI\_Config1. Si se grabara SPI, el
chip dejaría de responder por I\textsuperscript{2}C en funcionamiento normal. La hoja de datos indica
que tras el POR el chip atiende siempre por I\textsuperscript{2}C durante la ventana del modo comando,
de modo que sería recuperable, pero es mejor no correr el riesgo.

\section{Lecturas de la célula siempre ``antiguas''}

\paragraph{Síntoma.} En una versión del firmware no se registraba ni una sola muestra válida de la
célula.

\paragraph{Causa.} La lectura se hacía en dos fases, separadas \SI{2}{\milli\second}, quedándose con
la segunda. En modo de conversión continua leer no dispara una conversión, y el chip no había
producido un dato nuevo en \SI{2}{\milli\second}, así que la segunda lectura venía siempre marcada como
repetida. Una versión anterior funcionaba por casualidad: sus dos lecturas caían en vueltas distintas
del lazo, separadas \SI{50}{\milli\second}.

\paragraph{Solución.} Una única lectura por vuelta, aceptada sólo si el dato es nuevo. El ritmo
efectivo es el menor entre el del lazo y el de conversión del chip.

\section{Discrepancias con la hoja de datos detectadas durante la redacción}
\label{sec:prob-estado}

Al contrastar el firmware y la interfaz con la hoja de datos del ZSC31014 para esta memoria, se
encontraron dos discrepancias que no afectan a los ensayos presentados, pero que conviene corregir:

\begin{enumerate}
  \item \textbf{Bits de estado.} Según la hoja de datos, 00 = dato válido, \textbf{01 = modo comando},
        \textbf{10 = dato repetido} y 11 = diagnóstico. El firmware define
        \codigo{LC\_STATUS\_STALE = 1} y \codigo{LC\_STATUS\_COMMAND = 2}, es decir, intercambiados. En la
        lectura normal sólo se acepta 00, así que no hay efecto. En la confirmación del modo comando,
        \codigo{zscResp} acepta un primer byte con estado 10 o igual a \codigo{0x5A}. La respuesta
        real es \codigo{0x5A}, cuyos bits altos son 01, y el firmware la acepta por la segunda
        condición; la primera condición es errónea pero inocua.
  \item \textbf{Codificación de la ganancia.} Según la tabla de B\_Config, el campo PreAmp\_Gain [6:4]
        no es lineal: 000 = ×1,5, 100 = ×3, 001 = ×6, 101 = ×12, 010 = ×24, 110 = ×48, 011 = ×96 y
        111 = ×192. El índice de ganancia corresponde a los tres bits en orden inverso. El selector de
        la interfaz supone un código lineal (1 = ×3, 2 = ×6\ldots), así que para los códigos 1 a 6 la
        ganancia grabada no es la mostrada: elegir ``×3'' graba ×6, y elegir ``×24'' graba ×3. Los
        códigos 0 (×1,5) y 7 (×192, el valor por defecto del firmware) coinciden.
        \verificar{comprobar qué código de ganancia estaba grabado durante los ensayos}
\end{enumerate}

\section{Segunda célula de carga en el mismo bus}
\label{sec:prob-nau}

Se evaluó añadir una segunda célula con una Load Cell 4 Click (NAU7802) en el mismo bus
I\textsuperscript{2}C (programa \fichero{test\_loadcell\_dual\_v2}). Surgieron dos problemas:

\begin{itemize}
  \item \textbf{Convivencia en el bus.} Con la lectura en dos fases, el ZSC31014 no toleraba
        transacciones ajenas entre la petición y la recogida del dato: un reinicio de bus dirigido al
        NAU7802 estropeaba la siguiente lectura del ZSC31014. De ahí nació el árbitro de
        \fichero{i2c\_bus.h}.
  \item \textbf{Ruido del NAU7802.} Las medidas sobre la placa, con ganancia 128 y en reposo, dieron
        unas 18.000 cuentas pico a pico con el regulador interno a \SI{4.5}{\volt} (en caída, porque la
        placa se alimenta a \SI{3.3}{\volt}); unas 3.500 con el regulador a \SI{3.0}{\volt} y
        calibración interna de cero; y unas 350 al activar el \emph{chopper} del amplificador, que
        equivalen a \SI{0.49}{\micro\volt} en la entrada y coinciden con las entradas cortocircuitadas.
        También se comprobó que esa Click no alimenta AVDD por fuera, así que el conversor debe usar el
        regulador interno; en caso contrario satura en $-8\,388\,608$.
\end{itemize}

La segunda célula se retiró del firmware final para simplificar el lazo y el bus. Los ensayos
analizados usan sólo la célula del cable.

\section{Otros problemas}

\begin{description}[style=nextline,leftmargin=1.5em]
  \item[Órdenes perdidas por los atajos de teclado] Un byte de CRC igual a una letra de atajo hacía que
        se perdiera la orden y se reiniciara el servo (sección~\ref{chap:comunicaciones}). Se resolvió
        aceptando los atajos sólo con el decodificador en reposo.
  \item[Frecuencia de muestreo imposible] La interfaz llegó a mostrar \SI{3813}{\hertz}. Tras un reinicio
        del microcontrolador su reloj volvía a cero, la ventana de medida no se podaba nunca y se
        contaban todas las muestras acumuladas. Se detecta el salto de reloj y se reinicia la ventana.
  \item[Desplazamientos fantasma] Tras la programación de la célula, que bloquea el lazo cerca de un
        segundo, la integración multiplicaba ese hueco por la última velocidad. Se anula la base de
        tiempo tras cada operación bloqueante.
  \item[Órdenes que ``no hacían nada''] Una acción mal escrita en la interfaz se ignoraba sin avisar. Se
        añadió una confirmación para toda orden, con el motivo del fallo cuando lo hay.
  \item[Revisión A de la placa] Conflictos de red en el bloque RS-485, entradas analógicas en pines sin
        ADC y amplificadores sin alimentar (tabla~\ref{tab:erratas}).
\end{description}

\section{Lecciones aprendidas}

\begin{itemize}
  \item \textbf{Leer la hoja de datos entera}, no sólo las tablas de registros: la ventana del modo
        comando, la firma de la EEPROM, el umbral de POR y la codificación no lineal de la ganancia están
        en secciones distintas y todas han resultado críticas.
  \item \textbf{Medir antes de suponer}: la alimentación parásita sólo se descubrió con un polímetro, y
        el retardo de la ventana, con un barrido.
  \item \textbf{Cada operación debe responder}: las confirmaciones de órdenes y el modelo
        pregunta-respuesta de la célula han convertido fallos silenciosos en diagnósticos claros.
  \item \textbf{La sencillez da determinismo}: un lazo de periodo fijo con dueño único de cada bus ha
        dado mejor resultado que repartir el trabajo entre dos núcleos.
  \item \textbf{Conservar los datos brutos}: ha permitido recalibrar y reanalizar todos los ensayos para
        esta memoria sin repetirlos.
\end{itemize}
